\section{Thermodynamics and Thermochemistry}

\subsection{Ctrl+F keywords: enthalpy, Hess, formation enthalpy, combustion, entropy, Gibbs, spontaneity, equilibrium constant}
Use for heat, work, \(\Delta H\), \(\Delta S\), \(\Delta G\), equilibrium constants from thermodynamics, and temperature-dependent spontaneity.

\subsection{Symbols and notation}
\referencetable{tab:thermo:symbols}{Thermodynamic symbols and notation.}
\begin{tabular}{@{}>{$}l<{$} p{10.8cm}@{}}
\toprule
\text{Symbol} & \text{Meaning and usage}\\
\midrule
U,\ H,\ S,\ G & Internal energy, enthalpy, entropy, and Gibbs free energy.\\
q,\ w & Heat entering the system and work performed on the system.\\
p,\ q_p & Pressure and heat measured at constant pressure.\\
p_{\mathrm{ext}},\ \Delta V & External pressure and volume change in pressure-volume work.\\
T,\ R & Absolute temperature and gas constant.\\
\nu_i & Coefficient of species \(i\) in the balanced reaction.\\
{}^\circ,\ \mathrm{rxn},\ \mathrm f,\ \mathrm{comb} & Standard state; reaction, formation, and combustion subscripts.\\
\mathrm{prod},\ \mathrm{react},\ k & Products, reactants, and reaction-equation term index in summations.\\
K,\ Q & Equilibrium constant and current reaction quotient.\\
D & Average bond enthalpy used for bond-breaking and bond-forming estimates.\\
\bottomrule
\end{tabular}

\subsection{Core formulas}

\subsubsection{First law, heat, and work}
\begin{equation}
\Delta U=q+w, \qquad w=-p_{\mathrm{ext}}\Delta V
\label{eq:thermo:first-law-work}
\end{equation}
Heat into the system is positive. Expansion work is negative for the system when \(\Delta V>0\).

\subsubsection{Constant-pressure heat and enthalpy}
\begin{equation}
H=U+pV, \qquad q_p=\Delta H
\label{eq:thermo:enthalpy-constant-pressure}
\end{equation}
At constant pressure, measured heat equals the enthalpy change. \(\Delta H<0\) is exothermic and \(\Delta H>0\) is endothermic.

\subsubsection{Reaction enthalpy from standard formation enthalpies}
\begin{equation}
\boxed{\Delta H^\circ_{\mathrm{rxn}}=
\sum_{\mathrm{prod}}\nu_i\Delta H^\circ_{\mathrm f,i}
-\sum_{\mathrm{react}}\nu_i\Delta H^\circ_{\mathrm f,i}}
\label{eq:thermo:reaction-enthalpy-formation}
\end{equation}
Use coefficients from the balanced equation. Elements in their standard state have \(\Delta H^\circ_{\mathrm f}=0\).

\subsubsection{Hess's law and thermochemical scaling}
\begin{equation}
\Delta H^\circ_{\mathrm{overall}}=\sum_k \Delta H^\circ_k, \qquad
\ce{2A -> 2B}: \Delta H=2\Delta H_{\ce{A -> B}}
\label{eq:thermo:hess-scaling}
\end{equation}
Reversing a reaction changes the sign. Multiplying a reaction multiplies \(\Delta H\), \(\Delta S\), and \(\Delta G\).

\subsubsection{Combustion enthalpy}
\begin{equation}
\Delta H^\circ_{\mathrm{comb}}
=\sum\nu\Delta H^\circ_{\mathrm f}(\ce{CO2},\ce{H2O},\ldots)
-\sum\nu\Delta H^\circ_{\mathrm f}(\text{fuel},\ce{O2})
\label{eq:thermo:combustion-enthalpy}
\end{equation}
Balance complete combustion to \(\ce{CO2}\) and \(\ce{H2O}\) first. \(\Delta H^\circ_{\mathrm f}(\ce{O2(g)})=0\).

\subsubsection{Reaction enthalpy from average bond enthalpies}
\begin{equation}
\Delta H_{\mathrm{rxn}}\approx
\sum D(\text{bonds broken})-\sum D(\text{bonds formed})
\label{eq:thermo:bond-enthalpy-estimate}
\end{equation}
\begin{equation}
\begin{aligned}
\Delta H_{\mathrm{rxn}}&\approx
-\sum \Delta H_{\mathrm{bond}}(\text{bonds broken})
+\sum \Delta H_{\mathrm{bond}}(\text{bonds formed})\\
&\hspace{1cm}\text{when tabulated bond-formation values are negative}
\end{aligned}
\label{eq:thermo:signed-bond-data}
\end{equation}
Bond enthalpy data are usually positive bond-breaking magnitudes; if the supplied table instead prints negative bond-formation values, use Equation~\eqref{eq:thermo:signed-bond-data}. Count bonds from the balanced molecular reaction before applying either formula.

\subsubsection{Entropy change from standard molar entropies}
\begin{equation}
\Delta S^\circ_{\mathrm{rxn}}=
\sum_{\mathrm{prod}}\nu_iS^\circ_i-\sum_{\mathrm{react}}\nu_iS^\circ_i
\label{eq:thermo:reaction-entropy}
\end{equation}
Use \(\mathrm{J\,mol^{-1}K^{-1}}\) for entropy and convert to \(\mathrm{kJ}\) if combining with \(\Delta H\).

\subsubsection{Gibbs free energy and spontaneity}
\begin{equation}
\Delta G=\Delta H-T\Delta S
\label{eq:thermo:gibbs}
\end{equation}
\(\Delta G<0\) is spontaneous, \(\Delta G=0\) is equilibrium, and \(\Delta G>0\) is nonspontaneous under the stated conditions.

\subsubsection{Standard free energy and equilibrium constant}
\begin{equation}
\Delta G^\circ=-RT\ln K, \qquad K=e^{-\Delta G^\circ/RT}
\label{eq:thermo:equilibrium-constant}
\end{equation}
\(K>1\) corresponds to \(\Delta G^\circ<0\). Use kelvin and consistent energy units.

\subsubsection{Nonstandard Gibbs free energy}
\begin{equation}
\Delta G=\Delta G^\circ+RT\ln Q
\label{eq:thermo:nonstandard-gibbs}
\end{equation}
The reaction direction under current conditions is determined by \(\Delta G\), not by \(\Delta G^\circ\) alone.

\subsubsection{Spontaneity from signs and temperature}
\referencetable{tab:thermo:spontaneity-signs}{Spontaneity predicted from signs of \(\Delta H\) and \(\Delta S\).}
\begin{tabular}{@{}ccc@{}}
\toprule
\(\Delta H\) & \(\Delta S\) & Spontaneity\\
\midrule
\(-\) & \(+\) & All \(T\)\\
\(+\) & \(-\) & No \(T\)\\
\(-\) & \(-\) & Low \(T\)\\
\(+\) & \(+\) & High \(T\)\\
\bottomrule
\end{tabular}
\begin{equation}
T_{\mathrm{switch}}=\frac{\Delta H}{\Delta S}\qquad(\Delta G=0)
\label{eq:thermo:switch-temperature}
\end{equation}

\subsection{Exam recipe: reaction enthalpy and Hess's law}
\textbf{Use when:} the target is \(\Delta H^\circ_{\mathrm{rxn}}\) and formation enthalpies, combustion enthalpies, or thermochemical equations are provided.

\textbf{Given / Find:} Given a balanced or balanceable reaction and thermochemical data; find \(\Delta H^\circ_{\mathrm{rxn}}\) for the target reaction as written.

\textbf{Required references:} Equation~\eqref{eq:thermo:reaction-enthalpy-formation}, Equation~\eqref{eq:thermo:hess-scaling}, and, for combustion data, Equation~\eqref{eq:thermo:combustion-enthalpy}.

\textbf{Steps:}
\begin{enumerate}
    \item Balance the target reaction and retain its requested direction and stoichiometric scale.
    \item If \(\Delta H^\circ_{\mathrm f}\) values are given, insert the balanced coefficients in Equation~\eqref{eq:thermo:reaction-enthalpy-formation}.
    \item If reaction equations are given, reverse or scale each equation until species cancel to the target; change or scale each \(\Delta H^\circ\) by Equation~\eqref{eq:thermo:hess-scaling}.
    \item Add the adjusted enthalpy terms and report the value per target reaction as written; if heat per gram is requested, divide the molar reaction enthalpy by the mass of the specified substance consumed according to the balanced reaction.
\end{enumerate}

\textbf{Checks:} Standard-state elements contribute \(0\) in formation-data calculations; complete combustion is normally exothermic, so \(\Delta H^\circ_{\mathrm{comb}}<0\).

\textbf{Common traps:} Omitting stoichiometric coefficients; reversing a reaction without reversing the sign; changing the target equation after calculating its enthalpy; reporting \(\mathrm{kJ\,mol^{-1}}\) when the question asks for \(\mathrm{kJ\,g^{-1}}\).

\subsection{Exam recipe: reaction enthalpy from bond enthalpies}
\textbf{Use when:} the problem supplies bond enthalpies or bond energies for reactant and product bonds instead of formation enthalpies.

\textbf{Given / Find:} Given a molecular reaction and bond enthalpies; find an estimated \(\Delta H_{\mathrm{rxn}}\).

\textbf{Required references:} Equations~\eqref{eq:thermo:bond-enthalpy-estimate} and \eqref{eq:thermo:signed-bond-data}.

\textbf{Steps:}
\begin{enumerate}
    \item Balance the molecular reaction and draw or read enough structure to count each bond type broken in reactants and formed in products.
    \item If the data are positive dissociation energies \(D\), multiply each broken- and formed-bond count by \(D\), then subtract the formed total from the broken total using Equation~\eqref{eq:thermo:bond-enthalpy-estimate}.
    \item If the data are printed as negative bond-formation enthalpies, reverse the sign for bonds broken and retain the printed sign for bonds formed using Equation~\eqref{eq:thermo:signed-bond-data}.
\end{enumerate}

\textbf{Checks:} Forming more or stronger bonds than are broken should tend to make \(\Delta H\) negative; count all bonds after balancing.

\textbf{Common traps:} Adding positive formed-bond dissociation energies instead of subtracting them; applying the positive-\(D\) formula directly to negative bond-formation entries.

\subsection{Exam recipe: Gibbs spontaneity from signs and temperature}
\textbf{Use when:} the question gives signs or numerical values of \(\Delta H\) and \(\Delta S\) and asks when the process is spontaneous.

\textbf{Given / Find:} Given \(\Delta H\), \(\Delta S\), and possibly \(T\); find the sign of \(\Delta G\), the spontaneous temperature range, or the threshold temperature.

\textbf{Required references:} Equation~\eqref{eq:thermo:gibbs}, Table~\ref{tab:thermo:spontaneity-signs}, and Equation~\eqref{eq:thermo:switch-temperature}.

\textbf{Steps:}
\begin{enumerate}
    \item Convert \(\Delta S\) to the same energy unit basis as \(\Delta H\), normally \(\mathrm{kJ\,mol^{-1}K^{-1}}\).
    \item For sign-only questions, select the matching row in Table~\ref{tab:thermo:spontaneity-signs}.
    \item For a stated temperature, calculate \(\Delta G\) with Equation~\eqref{eq:thermo:gibbs} and classify by its sign.
    \item If spontaneity can switch with temperature, calculate \(T_{\mathrm{switch}}\) from Equation~\eqref{eq:thermo:switch-temperature} and state the spontaneous side of the threshold.
\end{enumerate}

\textbf{Checks:} A computed threshold temperature must be in kelvin and physically meaningful; \(\Delta G=0\) at the threshold.

\textbf{Common traps:} Mixing \(\mathrm{J}\) and \(\mathrm{kJ}\); interpreting \(\Delta G>0\) as spontaneous; using a threshold for the sign combinations that never switch.

\subsection{Exam recipe: equilibrium constant from thermodynamic data}
\textbf{Use when:} the target is \(K\) and standard Gibbs data or standard enthalpy and entropy data are supplied.

\textbf{Given / Find:} Given \(T\) and either \(\Delta G^\circ_{\mathrm{rxn}}\) or data for \(\Delta H^\circ_{\mathrm{rxn}}\) and \(\Delta S^\circ_{\mathrm{rxn}}\); find the dimensionless equilibrium constant \(K\).

\textbf{Required references:} Equation~\eqref{eq:thermo:gibbs}, Equation~\eqref{eq:thermo:reaction-entropy}, Equation~\eqref{eq:thermo:reaction-enthalpy-formation}, and Equation~\eqref{eq:thermo:equilibrium-constant}.

\pythontool{\texttt{equilibrium\_constant(...)} efter \(\Delta G^\circ\) er fundet; brug \texttt{constant("R")} hvis du regner eksponenten manuelt}

\textbf{Steps:}
\begin{enumerate}
    \item Balance the target reaction and calculate any missing \(\Delta H^\circ_{\mathrm{rxn}}\) or \(\Delta S^\circ_{\mathrm{rxn}}\) using the balanced coefficients.
    \item If \(\Delta G^\circ_{\mathrm{rxn}}\) is not given, calculate it at the stated \(T\) from Equation~\eqref{eq:thermo:gibbs}.
    \item Convert \(\Delta G^\circ_{\mathrm{rxn}}\) to \(\mathrm{J\,mol^{-1}}\) when using \(R\) in \(\mathrm{J\,mol^{-1}K^{-1}}\), and convert \(T\) to kelvin.
    \item Calculate \(K=e^{-\Delta G^\circ/RT}\) from Equation~\eqref{eq:thermo:equilibrium-constant}.
\end{enumerate}

\textbf{Checks:} \(\Delta G^\circ<0\) must give \(K>1\), and \(\Delta G^\circ>0\) must give \(K<1\).

\textbf{Common traps:} Using Celsius; missing the minus sign in the exponent; mixing energy units; using \(Q\) when the request is the standard-state equilibrium constant.
